\section*{Bab \ref{ch:g5:data} --- Soal Cerita dan Diagram}

\begin{solution}{exo:g5:data:1}
Karcis: $32 \times 7 = 224$. Seluruhnya: $240 + 224 = 464$.
\end{solution}

\begin{solution}{exo:g5:data:2}
Per satu kg: $7.50 \div 3 = 2.50$. Lima kg:
$5 \times 2.50 = 12.50$.
\end{solution}

\begin{solution}{exo:g5:data:3}
$87 \div 6 = 14.5$: masing-masing membayar $14.50$.
\end{solution}

\begin{solution}{exo:g5:data:4}
Dibelanjakan: $13.90 + 8.50 = 22.40$. Sisanya: $25 - 22.40 = 2.60$.
\end{solution}

\begin{solution}{exo:g5:data:5}
Lebih dari $50$ mm: Januari ($60$), Maret ($55$), April ($70$), Mei
($80$). Mei dikurangi Juni: $80 - 30 = 50$ mm.
\end{solution}

\begin{solution}{exo:g5:data:6}
Minggu $2$: $6$ cm. Pertumbuhan $4$ cm: antara minggu $2$ dan $3$
($6 \to 10$), dan antara minggu $4$ dan $5$ ($15 \to 19$).
\end{solution}

\begin{solution}{exo:g5:data:7}
Empat batang setinggi $4$, $7$, $3$, $6$ (skala setiap $1$ atau
$2$ gol). Seluruhnya: $4 + 7 + 3 + 6 = 20$ gol.
\end{solution}

\begin{solution}{exo:g5:data:8}
Per satu di dalam paket: $3.20 \div 8 = 0.40$. Hematnya per yogurt:
$0.50 - 0.40 = 0.10$ (sepuluh sen).
\end{solution}

\begin{solution}{exo:g5:data:9}
Cara 1: hari Sabtu $148 \times 9.50 = 1\,406$; hari Minggu
$95 \times 9.50 = 902.50$; \omterm{ex:g1:subtraction:difference}{selisihnya} $503.50$.

Cara 2: $148 - 95 = 53$ karcis lebih banyak, masing-masing bernilai
$9.50$: $53 \times 9.50 = 503.50$. Jawabannya sama --- \omterm{ex:g1:subtraction:difference}{selisih}
pemasukannya adalah pemasukan dari \omterm{ex:g1:subtraction:difference}{selisihnya}.
\end{solution}

\begin{solution}{exo:g5:data:10}
Per orang: $300 \div 4 = 75$ g. Untuk sepuluh orang:
$10 \times 75 = 750$ g beras.
\end{solution}

\begin{solution}{exo:g5:data:11}
Banyak soal yang benar --- misalnya: ``Miko membeli $3$ jus seharga
$1.20$ masing-masing dan sebuah roti isi seharga $3.40$. Ia membayar
dengan voucer $12.50$. Berapa sisa voucernya?'' Penyelesaian: jus
$3 \times 1.20 = 3.60$; seluruh belanja $3.60 + 3.40 = 7$; \omterm{def:g3:division:remainder}{sisanya}
$12.50 - 7 = 5.50$.

\end{solution}
